Standard sun-exposure guidance for vitamin D prescribes evenly spaced, fixed-time sessions. Existing models hold skin type fixed and have not compared spacing and dose jointly; considering both may reduce the exposure needed to reach and hold sufficiency.
We built a three-state ODE model coupling serum 25(OH)D, UV-based pigmentation, and stratum corneum thickening, adapting the synthesis-rate model of Young et al. (2021), the baseline-dependent model of Datta et al. (2017), and the stratum-corneum thickening data of de Boer et al. (2025). For each Fitzpatrick Skin Type (FST) I–VI and session budget, we used differential evolution to jointly search dose amount per session (two-phase or exponentially declining) and spacing against a uniform-spacing control. The doses were under sub-erythemal (skin-reddening) dosing and the model targeted a 50 nmol/L sufficiency floor while reducing skin cancer risk.
Two front-loading effects emerged. The chosen schedules changed exposure needed at UV index 6 from about 8 to 5.5 minutes at FST I and about 32 to 21 minutes at FST VI (evenly spaced vs. front-loaded schedules respectively). Separately, holding the per-session dose profile free in both arms and isolating session timing, front-loaded schedules lowered the combined cancer-risk index against an evenly spaced, equal-session-count control in 20 of 22 conditions tested, cutting the intermittency-driven (basal cell/melanoma) index by a median of ~22% while the cumulative-dose (squamous cell) index rose by a median of ~8%. These comparisons hold efficacy fixed at the sufficiency floor; single-seed estimates, order-of-magnitude.
These predictions remain preliminary: the synthesis scale rests mainly on FST I–II data, no field dosimetry exists for FST VI, and the model omits dietary and behavioural determinants. Still, guidance that varies spacing and dose with UV index and skin type could make sufficiency reachable with less sun exposure.
